\section{The annular affine manifold and its coefficient system}
\label{sec:affine-kummer}

\subsection{The quartic degeneration}

Let $M=\ZZ^2$ be the lattice of the fan of $\PP(1,1,2)$ (the lattice
denoted $N$ in toric geometry), with primitive ray generators
\[
 v_X=(1,0),\qquad v_Y=(-1,-2),\qquad v_Z=(0,1).
\]
The torus of $\PP(1,1,2)$ has character lattice
$M^\vee=\operatorname{Hom}(M,\ZZ)$, and we write $(s_1,s_2)$ for the
standard coordinates of the plane $M^\vee\otimes\RR$, dual to those of $M$.
The anticanonical polygon $\Delta\subset M^\vee\otimes\RR$, with vertices
given in the coordinates $(s_1,s_2)$, and its polar
$\Delta^\vee\subset M\otimes\RR$ are
\begin{equation}\label{eq:polygons}
 \Delta=\operatorname{conv}\{(-1,-1),(-1,1),(3,-1)\},
 \qquad
 \Delta^\vee=\operatorname{conv}\{v_X,v_Z,v_Y\}.
\end{equation}
The edge $[v_Y,v_X]$ has lattice length two and midpoint
$m_*=(0,-1)\in M$.  This is the unique non-basic edge of $\Delta^\vee$.

Consider the marginal quartic pencil
\begin{equation}\label{eq:quartic-pencil}
 E_s=\{X^4+Y^4+Z^2+sX^2Y^2=0\}\subset\PP(1,1,2)
\end{equation}
and the toric-boundary degeneration
\begin{equation}\label{eq:toric-degeneration}
 XYZ+\eta\bigl(X^4+Y^4+Z^2+sX^2Y^2\bigr)=0.
\end{equation}
For $\eta\ne0$, division by $\eta$ and the substitution
$Z'=Z+XY/(2\eta)$ give the same pencil with effective parameter
$s-1/(4\eta^2)$.  At $\eta=0$ the fiber is the toric boundary.

\begin{proposition}\label{prop:coherent-degeneration}
The family \eqref{eq:toric-degeneration} is a relative anticanonical section
of the coherent Mumford degeneration associated with the star subdivision of
$\Delta$ at the origin.  Its logarithmic smoothing charts are semistable on
the canonical stack.  At the coarse $A_1$ point the smoothing is Kummer of
index two.
\end{proposition}

\begin{proof}
The lattice points of $\Delta$ index the anticanonical Cox monomials by
\[
 m\longmapsto
 X^{\langle m,v_X\rangle+1}
 Y^{\langle m,v_Y\rangle+1}
 Z^{\langle m,v_Z\rangle+1}.
\]
The five terms in \eqref{eq:toric-degeneration} correspond to the origin and
the four boundary points indexing $X^4,Y^4,Z^2,X^2Y^2$.  Assigning height
zero to the origin and height one to the boundary gives the star
subdivision, hence the stated coherent degeneration.

On $Y\ne0$, with $x=X/Y$ and $z=Z/Y^2$, the point $x=z=0$ is a node of
the toric boundary at a smooth point of $\PP(1,1,2)$, and near it the family
has equation
\[
 xz+\eta(z^2+x^4+sx^2+1)=0.
\]
The parenthesized term is a unit at the node, so an étale change of
coordinates gives $xz=\eta$.  At $[0:0:1]$, pass to the canonical stack chart
with residual involution $(x,y)\mapsto(-x,-y)$.  The equation is again
$xy=\eta$ up to a unit.  On the coarse chart set
$\xi_0=x^2$, $\xi_1=xy$, $\xi_2=y^2$; since $\xi_0\xi_2=\xi_1^2$,
eliminating $\xi_1$ gives $\xi_0\xi_2=\eta^2$ up to a unit.
\end{proof}

Let $u$ denote the primitive character along the length-two slab $\rho$.
The slab function of $\rho$ may be written $f_\rho=1+a_\rho z^u+z^{2u}$, with a
free coefficient $a_\rho$, the \emph{slab parameter}.  The associated curve is
compared with
\eqref{eq:quartic-pencil} by
\begin{equation}\label{eq:curve-mirror-map}
 s^2=\frac{4a_\rho^2}{a_\rho^2-4},
 \qquad a_\rho^2=\frac{4s^2}{s^2-4}.
\end{equation}
Indeed, eliminating $U$ from
$U+V+a_\rho V^{-1}+U^{-1}V^{-2}=0$ gives
$Y^2=(V^2+a_\rho)^2-4$, and rescaling by a square root of $a_\rho^2-4$ gives the
quartic.  Equation~\eqref{eq:curve-mirror-map} is a comparison of the
anticanonical curves.  At the two values $a_\rho=\pm2$ the slab function
$f_\rho=(1\pm z^u)^2$ is a square, the curve $Y^2=(V^2+a_\rho)^2-4$ acquires a
node, and \eqref{eq:curve-mirror-map} gives $s=\infty$.  We call these
values the \emph{cusps}; the term refers to degenerate values of the
parameter, as for the cusps of a modular curve, and not to a cuspidal
singularity.

\subsection{Exponent lattices and cofactor directions}
\label{subsec:exponent-lattices}

Let $\ZZ^2_{x,y}=\ZZ e_1\oplus\ZZ e_2$ be the lattice of bidegrees of Laurent
monomials in the variables of $W=x^4+y^4$.  We write its elements as pairs
$(a,b)$ and read $(a,b)$ as the exponent of $x^ay^b$.  Bidegrees count
boundary markings: for the GKT potentials and generators recalled in
Section~\ref{subsec:gkt-recollections}, the bidegree of the monomial attached
to a disk graph, a monomial $x^{k_1}y^{k_2}$ of a potential or the
Hamiltonian $x^{k_1+1}y^{k_2+1}$ of a critical generator $X_k$, is the pair
of numbers of boundary tails of twists $(2,0)$ and $(0,2)$ on the graph.  The
coordinate divisors give
\begin{equation}\label{eq:iota}
 \iota:\ZZ^2_{x,y}\longrightarrow M,
 \qquad
 \iota=\begin{pmatrix}1&-1\\0&-2\end{pmatrix},
 \qquad [M:\iota(\ZZ^2_{x,y})]=2,
\end{equation}
so that $\iota(e_1)=v_X$, $\iota(e_2)=v_Y$ and $\iota(e_1+e_2)=2m_*$.  The
support map for walls is the cofactor
\begin{equation}\label{eq:cofactor}
 L=\det(\iota)\iota^{-T}
  =\begin{pmatrix}-2&0\\1&1\end{pmatrix}.
\end{equation}
For a pair $(a,b)$, whether a bidegree, an element of $M$ in standard
coordinates or a vector in the coordinates $(s_1,s_2)$, write
$(a,b)^\perp=(b,-a)$.  Then, for $q\in\ZZ^2_{x,y}$,
\begin{equation}\label{eq:cofactor-identity}
 (Lq)^\perp=\iota(q^\perp),
 \qquad L^T\iota=\det(\iota)\,\mathrm{Id},
\end{equation}
where $(Lq)^\perp$ is formed in the coordinates $(s_1,s_2)$ and
$\iota(q^\perp)$ in the standard coordinates of $M$, which are dual to them.
By the second identity, $\iota(q^\perp)\in M$ annihilates $Lq$.
Section~\ref{subsec:cofactor-supports} coorients the wall of direction $Lq$
by this normal.

The Landau--Ginzburg variables enter through the linear map
\begin{equation}\label{eq:tilde-M}
 F=\frac14\iota,
 \qquad \widetilde M=F(\ZZ^2_{x,y})\supset M,
 \qquad [\widetilde M:M]=8,
\end{equation}
and the dictionary $x^ay^b\mapsto z^{F(a,b)}$.  The map $F$ is the unique
linear map sending the exponents $(4,0)$ and $(0,4)$ of the two monomials of
$W$ to the rays $v_X$ and $v_Y$ of the coordinate divisors $\{X=0\}$ and
$\{Y=0\}$; the factor $\frac14$ is the reciprocal of the degree of $W$, that
is, the common weight of $x$ and $y$.  The dictionary sends
$x^4$, $y^4$, $x^2y^2$ and $xy$ to $z^{v_X}$, $z^{v_Y}$, $z^{m_*}$ and
$z^{F(1,1)}$.  Thus $\widetilde M$ is, under the dictionary, the lattice of
exponents of Laurent monomials in $x$ and $y$, while the fan lattice $M$ is the character lattice
of the dual torus $\Spec\Bbbk[M]$, on which the Laurent polynomial
$U+V+a_\rho V^{-1}+U^{-1}V^{-2}$ above and the slab functions live.  As for the
Landau--Ginzburg mirror of a toric variety, whose monomials have the rays of
the fan as exponents, the exponents of slab functions, wall functions and
derivations lie on the Legendre-dual side, in $M$ and $\widetilde M$, while
\eqref{eq:toric-degeneration} is written on the torus of $\PP(1,1,2)$, with
character lattice $M^\vee$.  The character
of $x^ay^b$ lies in $M$ exactly when $b$ is even and $a\equiv b\pmod4$.  In
particular $4\widetilde M=\iota(\ZZ^2_{x,y})$ has index two in $M$.

\emph{The dictionary.}  The Landau--Ginzburg variables $x,y$ are the Cox
coordinates $X,Y$ of $\PP(1,1,2)$, of weight one, and the Cox coordinate $Z$,
of weight two, is the suspension variable: $W(X,Y)+Z^2$ is the quartic
\eqref{eq:quartic-pencil} at $s=0$, and the marginal monomial $x^2y^2$ is the
deformation direction $X^2Y^2$ of the pencil.  The dictionary therefore
matches the pure powers $x^4$ and $y^4$ of the two Landau--Ginzburg variables
with the rays $v_X$ and $v_Y$ of the divisors $\{X=0\}$ and $\{Y=0\}$ of the
corresponding Cox coordinates.  The third ray $v_Z$, that of the divisor
$\{Z=0\}$ of the suspension variable, corresponds to no monomial of $W$.  The
extension from $x^4,y^4$ to all monomials is linear because the dictionary
is to be multiplicative: it must send products of monomials to products of
characters, so that it extends to a homomorphism from Laurent polynomials in
$x,y$ to $\Bbbk[\widetilde M]$ and carries derivations to derivations
(Proposition~\ref{prop:integral-cofactor}).  A multiplicative assignment of
characters to monomials is a homomorphism of exponent lattices, and it is
determined by its values on $(4,0)$ and $(0,4)$, which span
$\ZZ^2_{x,y}\otimes\QQ$.  On the side of the CPS complex, the characters of
$M$ are the exponents of the slab functions and of the Laurent polynomial
$U+V+a_\rho V^{-1}+U^{-1}V^{-2}$ above.  In the coordinates $U=z^{v_X}$ and
$V=z^{v_Z}$ of the torus $\Spec\Bbbk[M]$, its monomials $U$,
$U^{-1}V^{-2}$ and $V$ have the three rays $v_X$, $v_Y$ and $v_Z$ of the fan
as exponents, and $V^{-1}$ has the midpoint $m_*$ of the non-basic edge.  The
dictionary sends $x^4$, $y^4$ and $x^2y^2$ to $U$, $U^{-1}V^{-2}$ and
$V^{-1}$.  The coefficient of $V^{-1}$ is the slab parameter $a_\rho$, which is
related to the marginal parameter $s$ by \eqref{eq:curve-mirror-map}, and the
remaining monomial $V$ has as exponent the ray $v_Z$ of the divisor of the
suspension variable.

\emph{Coordinates.}  The map $F$ is an isomorphism of
$\ZZ^2_{x,y}$ onto $\widetilde M$, and we give $\widetilde M$ the coordinates
it transports.  The \emph{$F$-coordinates} of $m\in\widetilde M$ are the
coordinates of $F^{-1}(m)\in \ZZ^2_{x,y}$, that is, the coordinates of $m$ in
the basis $F(1,0),F(0,1)$; the character of $x^ay^b$ has $F$-coordinates
$(a,b)$, and we write $z^{(a,b)}$ for $z^{F(a,b)}$.  Dually, the
$F$-coordinates of $n\in M^\vee\otimes\QQ$, in particular of
$n\in\widetilde M^\vee$, are the coordinates of $F^Tn$, that is, the
coordinates of $n$ in the dual basis; the element with $F$-coordinates
$(c,d)$ is $F^{-T}(c,d)$.  Throughout the paper, exponents of characters of
$\widetilde M$, matrices acting on $\widetilde M$, and determinants of pairs
of elements of $\widetilde M$ are written in $F$-coordinates.  In particular
$\kappa=(1,1)$ is the exponent of $xy$, and with $z_1=z^{(1,0)}$ and
$z_2=z^{(0,1)}$, the characters of $x$ and $y$, the form
$z^\kappa\,d\log z_1\wedge d\log z_2$ of Section~\ref{sec:jacobi-theta}
corresponds to $dx\wedge dy$.  Standard coordinates of
$M\otimes\QQ$ are used only for elements of $M$
and for values such as $F(2,1)=(\frac14,-\frac12)$.  In passages that use
them, an element of $\widetilde M$ with $F$-coordinates $(a,b)$ is written
$F(a,b)$, never as a bare pair.  Pairs describing points and directions in
the plane $M^\vee\otimes\RR$ of $\Delta$ are in the coordinates $(s_1,s_2)$
of \eqref{eq:polygons} and are marked as such, as are $F$-coordinates of
elements of $M^\vee\otimes\QQ$.

Let $q=(q_1,q_2)\in\ZZ^2_{x,y}$, so that $q^\perp=(q_2,-q_1)$, and denote by
$\partial_v$ the invariant logarithmic derivation in direction $v\in\ZZ^2$,
so that $\partial_v(x^ay^b)=\langle v,(a,b)\rangle x^ay^b$.  The derivation
\[
 X^\kappa_q=x^{q_1-1}y^{q_2-1}\,\partial_{q^\perp}
\]
is the Hamiltonian vector field of $x^{q_1}y^{q_2}$ for $dx\wedge dy$, in the
sense that $(dx\wedge dy)(X^\kappa_q,\,\cdot\,)=d(x^{q_1}y^{q_2})$; the
superscript records the shift of the exponent by $\kappa=(1,1)$, the
exponent of $xy$.  Put
\begin{equation}\label{eq:normal-q}
 \widetilde n_q=F^{-T}q^\perp\in\widetilde M^\vee,
\end{equation}
where $F^{-T}$ is computed in the standard coordinates of
$M^\vee\otimes\QQ=\operatorname{Hom}(M,\QQ)$; thus
$\widetilde n_q$ is the element with $F$-coordinates $q^\perp$.
Consequently, in $F$-coordinates, where $\partial_v$ denotes the derivation
of $\Bbbk[\widetilde M]$ in the direction $F^{-T}v$, the transported
derivation \eqref{eq:transported-derivation} below is
$z^{q-\kappa}\partial_{q^\perp}$.  This is the field $X^\kappa_q$ of
\eqref{eq:shifted-Hamiltonian} and of the introduction, and we use the same
notation for both.

\begin{proposition}\label{prop:integral-cofactor}
Under the dictionary $x^ay^b\mapsto z^{F(a,b)}$, every logarithmic
Hamiltonian derivation $X^\kappa_q$ becomes an integral monomial derivation
of $\CC[\widetilde M]$.  More precisely,
\begin{equation}\label{eq:transported-derivation}
 F_*X^\kappa_q(z^m)
 =\langle\widetilde n_q,m\rangle z^{m+F(q-\kappa)},
 \qquad m\in\widetilde M.
\end{equation}
The derivation annihilates the wall character $z^{F(q)}$ and shifts
exponents by $F(q-\kappa)$.
\end{proposition}

\begin{proof}
For $q'\in\ZZ^2_{x,y}$,
\[
 \langle F^{-T}q^\perp,Fq'\rangle=\langle q^\perp,q'\rangle.
\]
This proves integrality and \eqref{eq:transported-derivation}.  Since
$\langle\widetilde n_q,F(q)\rangle=\langle q^\perp,q\rangle=0$, the
derivation annihilates $z^{F(q)}$.
\end{proof}

\emph{Wall directions.}  In the scattering diagrams of
\cite{grosspandharipandesiebert2010,grossetal2016theta} exponents and wall
supports lie in one real plane, and a wall is tangent to the exponent of its
function.  Here the two lie on opposite sides of the discrete Legendre
transform (Section~\ref{subsec:parity-cover}): the character $F(q)$ of the
Hamiltonian $x^{q_1}y^{q_2}$ of $X^\kappa_q$ lies in
$\widetilde M\subset M\otimes\QQ$, on the Legendre-dual side, whereas cuts
and wall supports are drawn in the plane $M^\vee\otimes\RR$ of $\Delta$,
through its central vertex $O$, the origin.  A support rule must therefore
carry exponents across the transform.  The transform matches each cell with a
dual cell whose tangent space is the annihilator of its own: the segment from
$O$ to a vertex $p$ of $\Delta$ (a \emph{star edge}) is matched with the edge
$\Delta^\vee\cap\{\langle p,\cdot\rangle=-1\}$ of the polar polygon, which
is parallel to the annihilator of $p$.  The singular points $A,B,C$ of
Section~\ref{subsec:parity-cover} are the midpoints of the star edges, $B$ on
the edge to $(-1,1)$, and each star edge is annihilated by the invariant
character $\bar u_V$ of the monodromy about its singular point $V$
(Appendix~\ref{app:calculations}).  If a support through $O$ is regarded as
an edge of a refinement of the star subdivision, the same rule matches a wall
tangent to $F(q)$ on the Legendre-dual side with a support along the
annihilator of $F(q)$, that is, along $\widetilde n_q=F^{-T}q^\perp$, the
direction of the derivation $X^\kappa_q$.  We do not use this rule, because it does not
respect the exchange of $x$ and $y$.  The exchange $\sigma(a,b)=(b,a)$ acts
on $M$ by $S=F\sigma F^{-1}$, which fixes $v_Z$ and exchanges $v_X$ and $v_Y$,
and on the plane of $\Delta$ by $S^T$, an involution of $\Delta$ that fixes
the star edge through $B$ pointwise and exchanges $A$ and $C$.  Since
$(\sigma q)^\perp=-\sigma(q^\perp)$, the rule above satisfies
$\widetilde n_{\sigma q}=-S^T\widetilde n_q$: the involution carries the
support of $X^\kappa_q$ to the opposite of the support of $X^\kappa_{\sigma q}$.
No choice of orientations repairs this: $\sigma$ fixes $(1,1)$, while $S^T$
reverses both rays of the annihilator of $F(1,1)$, the line $s_2=0$, which
moreover misses $A$, $B$ and $C$.

We place the wall of $X^\kappa_q$ along the cofactor direction
$Lq\in M^\vee\otimes\RR$ instead.  By \eqref{eq:cofactor},
$L=-\frac12F^{-T}$, so that
\begin{equation}\label{eq:cofactor-pairing}
 \langle Lq,F(q')\rangle=-\tfrac12\,q\cdot q',
 \qquad q,q'\in\ZZ^2_{x,y},
\end{equation}
with $q\cdot q'$ the standard scalar product of bidegrees, and $L\sigma=S^TL$.
Equivalently,
\[
 Lq=-8(\iota\iota^T)^{-1}F(q)=d\psi\bigl(F(q)\bigr),
 \qquad
 \psi\bigl(F(a,b)\bigr)=-\tfrac14\bigl(a^2+b^2\bigr),
\]
where $d\psi:M\otimes\RR\to M^\vee\otimes\RR$ is the gradient map, or
Legendre transform, of the quadratic function $\psi$ on $M\otimes\RR$.  Thus
$d\psi$ carries the ray of each exponent to the support of its wall.  Here
$\psi$ is an auxiliary quadratic function, unrelated to the polarization
$\varphi$; in $F$-coordinates on $\widetilde M$ and the dual $F$-coordinates
on $\widetilde M^\vee$ the rule is $q\mapsto-\frac12q$.  What singles out
rules of this kind is the exchange: a linear map
$P:\ZZ^2_{x,y}\otimes\RR\to M^\vee\otimes\RR$ satisfies $P\sigma=S^TP$
exactly when $P=F^{-T}Q$ for a matrix $Q$ commuting with $\sigma$.  Such a $Q$
is symmetric, so $P$ is the composite of $F$ with the gradient map of the
$\sigma$-invariant quadratic function $F(v)\mapsto\frac12v^TQv$, and $P$
sends $(1,1)$ into the line through $O$ and $B$.  The cofactor is the case
$Q=-\frac12\mathrm{Id}$.  Every negative definite $Q$ gives supports in the
cyclic order of the exponents, with the diagonal support running from $O$
along the star edge to $B$, and for all $Q$ except the multiples of
$-\bigl(\begin{smallmatrix}7&-5\\-5&7\end{smallmatrix}\bigr)$, whose chords pass
through $A$ and $C$, the supports avoid the singular points; among these the
choice is a convention.  The map $d\psi$ is symmetric and negative definite,
of determinant $16$ for the standard orientation of $M$ and the dual
orientation of $M^\vee$, so the supports occur about $O$ in the same cyclic
order as the characters $F(q)$ about the origin.  This cyclic order is the only property of the
placement used in the proof of consistency at the joints
(Theorem~\ref{thm:annular-consistency}): a path passing above $O$ from left
to right meets the supports in order of increasing slope of $q$, and its
path-ordered product is the ordered factorization
\eqref{eq:singleton-increasing-slope}.

The pairs $(2,1)$ and $(1,2)$ are the exponents of $x^2y$ and $xy^2$ in
$\ZZ^2_{x,y}$.  Their characters $F(2,1)=(\frac14,-\frac12)$ and
$F(1,2)=(-\frac14,-1)$ do not lie in $M$, and the corresponding wall supports
have the cofactor directions $L(2,1)=(-4,3)$ and $L(1,2)=(-2,3)$ in the
coordinates $(s_1,s_2)$.

The lattice $M$ is generated by $F(4,0)=v_X$ and $F(-2,-2)=v_Z$.  The
matrix of their $F$-coordinates $(4,0)$ and $(-2,-2)$ has determinant $-8$
and entries with greatest common divisor two, so the quotient is
\begin{equation}\label{eq:sector-group}
 \widetilde M/M\simeq\ZZ/2\oplus\ZZ/4.
\end{equation}
We call $\widetilde M/M$ the \emph{fractional-exponent group}: it records an exponent in
$\widetilde M$ modulo the integral exponents $M$.  The classes
$\alpha=[F(1,0)]$ and $\beta=[F(0,1)]$, by which the singleton fields
$X^\kappa_{(2,1)}$ and $X^\kappa_{(1,2)}$ shift exponents, have order four,
whereas $\alpha+\beta$ has order two.  The total degree modulo four,
\begin{equation}\label{eq:diagonal-charge}
 \deg_4:\widetilde M/M\longrightarrow\ZZ/4,
 \qquad \deg_4([F(a,b)])=a+b\pmod4,
\end{equation}
is a homomorphism with kernel generated by $\alpha-\beta$.

\subsection{The parity cover}
\label{subsec:parity-cover}

We keep the two sides of the discrete Legendre transform apart.  The polygon
$\Delta$ lies in $M^\vee\otimes\RR$.  With its star subdivision at
the origin $O$, the \emph{central vertex}, and the singular points $A,B,C$ at
the midpoints of the three
star edges, it carries the integral-affine structure with straight boundary
of Appendix~\ref{app:calculations}, whose tangent lattice is
$M^\vee$.  Cuts, wall supports, flares and developed pictures are
drawn in $\Delta$.  The dual lattice $M$ is the tangent lattice of the
Legendre-dual side.  Exponents of slab
functions, wall functions and derivations lie in $\widetilde M\subset
M\otimes\QQ$, and the polarization $\varphi$ is a piecewise-linear function
on the Legendre-dual side whose kinks are integral vectors of
$M^\vee$.  As topological manifolds we identify the Legendre-dual
complex with the interior of $\Delta$ by the correspondence between cells and
dual cells of the discrete Legendre transform
\cite[\S1.4]{grosssiebert2006}, under which the singular points correspond
to $A,B,C$; in particular $\varphi$ is a function on the interior of
$\Delta$, and it pulls back along the covers of $\Delta$ constructed below.
Along a loop, the monodromy of $M$ is the inverse
transpose of the monodromy of $M^\vee$.  The \emph{width} of a
unipotent monodromy $T\neq\mathrm{Id}$ is the positive integer $w$ with
$T-\mathrm{Id}=\pm w\,u\,n^T$ for primitive vectors $u$ and $n$, so that $T$
is conjugate in $\mathrm{GL}_2(\ZZ)$ to
$\left(\begin{smallmatrix}1&-w\\0&1\end{smallmatrix}\right)$; $T$ is
\emph{primitive} when $w=1$.  Widths, traces and
primitivity of a monodromy do not change under inverse transposition, so
statements about them refer to either lattice.  On the double cover
constructed below, $\widetilde M$ and its dual
$\widetilde M^\vee=F^{-T}(\ZZ^2)\subset M^\vee\otimes\QQ$ take the
places of $M$ and $M^\vee$; we call the pulled-back structure with
these lattices the \emph{$\widetilde M$-affine structure}.

For the loops $\gamma_A,\gamma_B,\gamma_C$ oriented as in
Appendix~\ref{app:calculations}, the monodromies of $M$ about the three
singular points are, in standard coordinates,
\begin{equation}\label{eq:disk-monodromies}
 T_A=\begin{pmatrix}0&-1\\1&2\end{pmatrix},\qquad
 T_B=\begin{pmatrix}3&-2\\2&-1\end{pmatrix},\qquad
 T_C=\begin{pmatrix}4&-1\\9&-2\end{pmatrix}.
\end{equation}
Only $T_B$ preserves $\widetilde M$.  Indeed, conjugating the other two by
$F$ gives half-integral matrices; see Appendix~\ref{app:calculations}.  Let
$\Delta_{\mathrm{aff}}^\circ$ be the affine base disk with the three
singular points $A,B,C$ removed, and define
\begin{equation}\label{eq:parity-character}
 \chi(\gamma_A)=\chi(\gamma_C)=1,
 \qquad \chi(\gamma_B)=0
 \quad\text{in }\ZZ/2.
\end{equation}

\begin{theorem}[Annular parity cover]\label{thm:parity-cover}
The connected double cover of $\Delta_{\mathrm{aff}}^\circ$ classified by
$\chi$ extends over the compact disk as a cover
$p:\overline\Acal\to\Delta$ ramified at $A$ and $C$, and $\overline\Acal$ is
an annulus.  The point $B$ has two unramified lifts $B_0,B_1$, which are
primitive focus--focus singularities of the $\widetilde M$-affine structure.
The ramification points $\widetilde A,\widetilde C$ over $A,C$ are singular
points of a different kind: near each of them the structure is, as a real
affine structure, the pullback of the focus--focus germ at $A$, respectively
$C$, under the branched double cover $z\mapsto z^2$; the lattice is
$\widetilde M^\vee$, which the monodromy of a single loop about $A$ or $C$
does not preserve.  Their linear monodromies on $\widetilde M$ are
primitive unipotent, but a small loop about $\widetilde A$ or $\widetilde C$,
based on its invariant line, develops onto a closed curve winding twice around the developed singular
point, whereas at a focus--focus point it winds once.  The pullback
polarization becomes integral after multiplication by four, and
$4p^*\varphi$ is the minimal integral multiple.
\end{theorem}

\begin{proof}
Cut the disk along an arc from $A$ to $C$ disjoint from $B$, and glue two
copies crosswise along it.  We take for this arc the \emph{branch segment},
the straight segment from $A$ to $C$ in the coordinates $(s_1,s_2)$.  The
local model at $A$ and at $C$ is $z\mapsto z^2$.  Since the boundary loop has parity $1+0+1=0$, it has two lifts.
Writing $\chi_{\mathrm{top}}$ for the Euler characteristic, Riemann--Hurwitz
gives $\chi_{\mathrm{top}}(\overline\Acal)=2\chi_{\mathrm{top}}(\Delta)-2=0$,
since the disk $\Delta$ has Euler characteristic one; a connected oriented
surface with Euler characteristic zero and two boundary components is an
annulus.

Appendix~\ref{app:calculations} computes the monodromy of the parity
subgroup on $\widetilde M$.  With the base point in the sheet $S_0$ of
Section~\ref{subsec:kummer-charts}, the monodromies about the four singular
points are, in $F$-coordinates,
\begin{equation}\label{eq:lifted-monodromies}
\begin{aligned}
 H_A&=\begin{pmatrix}-2&9\\-1&4\end{pmatrix},&
 H_{B,0}&=\begin{pmatrix}2&1\\-1&0\end{pmatrix},\\
 H_{B,1}&=\begin{pmatrix}-14&25\\-9&16\end{pmatrix},&
 H_C&=\begin{pmatrix}-2&1\\-9&4\end{pmatrix}.
\end{aligned}
\end{equation}
Their rank-one parts factor as
\begin{equation}\label{eq:primitive-factorizations}
\begin{aligned}
 H_A-\mathrm{Id}&=\binom31(-1,3),&
 H_{B,0}-\mathrm{Id}&=\binom{1}{-1}(1,1),\\
 H_{B,1}-\mathrm{Id}&=\binom53(-3,5),&
 H_C-\mathrm{Id}&=\binom13(-3,1),
\end{aligned}
\end{equation}
and in each factorization both vectors are primitive.  Hence every
$H_V$ is primitive unipotent on $\widetilde M$.  The points $B_0,B_1$ are
unramified lifts of $B$, so their germs are copies of the germ at $B$, a
singular point of the disk with unipotent monodromy of width two on $M$; with
respect to $\widetilde M$ the monodromies $H_{B,0},H_{B,1}$ are primitive.  A
small loop about $\widetilde A$ maps to the loop $\gamma_A$ traversed twice.
Develop it from a point of the invariant line through $\widetilde A$.  The
image is the developed image of $\gamma_A$ followed by its image under the
affine monodromy about $A$, which fixes the invariant line pointwise and
preserves orientation; each is a closed curve winding once around the
developed image of $A$, so the lift winds twice.  The same argument applies
at $\widetilde C$.  Winding numbers of developed loops are invariant under
affine isomorphism of germs, so the germs at $\widetilde A,\widetilde C$ are
not focus--focus germs.  The preimage of the star edge through $A$ consists
of two arcs crossing at $\widetilde A$, so $\widetilde A$ and $\widetilde C$
are four-valent vertices of the lifted decomposition.

The kinks of $\varphi$ across the three bounded edges of the Legendre-dual
decomposition are, up to sign, the vertices $p_A,p_B,p_C$ of $\Delta$,
since the star edges from $O$ to these vertices have lattice length one.
Their $F$-coordinates have denominator exactly four
(Appendix~\ref{app:calculations}).  The kinks across the unbounded edges are
the edge vectors of $\Delta$, whose $F$-coordinates are integral.
Multiplication by four therefore makes $p^*\varphi$ integral, and no
smaller positive multiple does.
\end{proof}

The two boundary circles of $\overline\Acal$ are the lifts of the boundary
loop $\gamma_A\gamma_B\gamma_C$ of $\Delta$, whose monodromy on $M$ is, in
standard coordinates,
$T_AT_BT_C=\begin{pmatrix}1&0\\-8&1\end{pmatrix}$, of width eight.  Their
monodromies $U_0,U_1$ on $\widetilde M$ are computed in
Appendix~\ref{app:calculations} and have normal forms
\begin{equation}\label{eq:boundary-holonomies}
 \begin{pmatrix}1&-4\\0&1\end{pmatrix},
 \qquad
 \begin{pmatrix}1&-16\\0&1\end{pmatrix}.
\end{equation}
Their widths are therefore $4$ and $16$.

Let $\Acal$ be the interior of $\overline\Acal$, and let $\Acal^\circ$ be the
complement in $\Acal$ of the four singular points
$\widetilde A,B_0,B_1,\widetilde C$.  Loops in $\Acal^\circ$ that are freely
homotopic in $\overline\Acal$ to a core circle of the annulus are not all
freely homotopic in $\Acal^\circ$; they differ by loops about the singular
points, and their holonomies need not have the same trace.  We fix the
\emph{core loop} $\gamma_{\mathrm{core}}=\gamma_C\gamma_A^{-1}$, based in the
sheet $S_0$ of Section~\ref{subsec:kummer-charts}.  Write
$\gamma_{\widetilde A}=\gamma_A^2$, $\gamma_{B_0}=\gamma_B$,
$\gamma_{B_1}=\gamma_A\gamma_B\gamma_A^{-1}$ and
$\gamma_{\widetilde C}=\gamma_C^2$ for the loops about the four singular
points.  In the fundamental group the two boundary loops factor as
\[
 \gamma_A\gamma_B\gamma_C
 =\gamma_{B_1}\gamma_{\mathrm{core}}^{-1}\gamma_{\widetilde C},
 \qquad
 \gamma_A^2\gamma_B\gamma_C\gamma_A^{-1}
 =\gamma_{\widetilde A}\gamma_{B_0}\gamma_{\mathrm{core}},
\]
so $\gamma_{\mathrm{core}}$ is the core loop separating
$\{\widetilde A,B_0\}$ from $\{B_1,\widetilde C\}$.  Its linear holonomy is
the hyperbolic matrix $K$ of \eqref{eq:intro-K}, of trace $18$, and
$U_0=H_{B,1}K^{-1}H_C$, $U_1=H_AH_{B,0}K$.  Other loops of this kind
have other traces: for instance $\gamma_A\gamma_C$ has holonomy
$K^{-1}H_C$, of trace $-14$, which exchanges $\alpha$ and $\beta$.

\begin{figure}[t]
\centering
\vspace{0.5\baselineskip}
\begin{adjustbox}{max width=\textwidth,center}
\begin{tikzpicture}[x=0.82cm,y=0.82cm,
  every node/.style={font=\small},
  focus/.style={circle,fill=black,inner sep=1.5pt},
  charge/.style={-{Latex[length=2.2mm]},line width=0.75pt},
  affine/.style={line width=0.8pt}]
  \path[use as bounding box] (-3.05,-2.05) rectangle (11.65,2.72);
  % Left panel: polygon and cofactor directions.
  \begin{scope}[shift={(0,0)}]
    \node[font=\bfseries] at (1,2.48) {(a) Cofactor geometry};
    \fill[blue!5] (-1,-1) -- (-1,1) -- (3,-1) -- cycle;
    \draw[affine] (-1,-1) -- (-1,1) -- (3,-1) -- cycle;
    \coordinate (pA) at (-1,-1);
    \coordinate (pB) at (-1,1);
    \coordinate (pC) at (3,-1);
    \coordinate (O) at (0,0);
    \fill (pA) circle (1.2pt) node[below left=1pt] {$p_A$};
    \fill (pB) circle (1.2pt) node[above left=1pt] {$p_B$};
    \fill (pC) circle (1.2pt) node[below right=1pt] {$p_C$};
    \fill (O) circle (1.2pt) node[below right=1pt] {$O$};
    \draw[orange!85!black,line width=1.15pt] (O) -- (pB);
    % legend for the highlighted star edge through the width-two point B
    \draw[orange!85!black,line width=1.15pt] (0.0,-0.78) -- (0.4,-0.78)
      node[right=1pt,text=orange!70!black] {width $2$};
    \draw[charge,blue!75!black] (O) -- ++(-1.72,1.29)
      node[above left=-1pt] {$L(2,1)=(-4,3)$};
    \draw[charge,teal!70!black] (O) -- ++(-1.08,1.62)
      node[above right=-2pt] {$L(1,2)=(-2,3)$};
    \node[align=center,text width=3.4cm] at (1.64,1.12)
      {$\Delta=\operatorname{conv}\{p_A,p_B,p_C\}$\\[-1pt]
       $L=\det(\iota)\iota^{-T}$};
  \end{scope}

  % Separator.
  \draw[gray!45] (4.15,-1.35) -- (4.15,2.35);

  % Right panel: annular parity cover; the dashed arc is a lift of the branch segment.
  \begin{scope}[shift={(7.65,0.45)}]
    \node[font=\bfseries] at (0,2.38) {(b) Parity cover};
    \fill[blue!5,even odd rule]
      (0,0) circle (1.62cm)
      (0,0) circle (0.68cm);
    \draw[affine] (0,0) circle (1.62cm);
    \draw[affine] (0,0) circle (0.68cm);
    \node[focus,label={[xshift=-1pt,yshift=1pt]above left:$\widetilde A$}] at (-0.95,0.98) {};
    \node[focus,label={[xshift=1pt,yshift=1pt]above right:$\widetilde C$}] at (0.95,0.98) {};
    \node[focus,label=below left:$B_0$] at (-0.62,-0.60) {};
    \node[focus,label=below right:$B_1$] at (0.62,-0.60) {};
    \node[blue!65!black,inner sep=1pt] at (0,-2.2) {width $4$};
    \node[blue!65!black,inner sep=1pt,font=\scriptsize] at (0,0.02) {width $16$};
    \draw[dashed,gray!75,line width=.7pt] (-0.95,0.98)
      .. controls (0,0.50) .. (0.95,0.98);
    \node[align=center,text width=2.45cm,anchor=west] at (1.85,0.12)
      {$\chi(\gamma_A)=\chi(\gamma_C)=1$\\
       $\chi(\gamma_B)=0$};
  \end{scope}
\end{tikzpicture}

\end{adjustbox}
\caption{The quartic polygon and its annular parity cover.  (a) The polygon
$\Delta$ with vertices $p_A,p_B,p_C$, in the coordinates $(s_1,s_2)$; the
highlighted star edge from $O$ to $p_B$ contains the singular point $B$,
whose monodromy has width two.  The first two singleton supports are the
cofactor directions $L(2,1)$ and $L(1,2)$.  (b) Two copies of the disk, cut
along the branch segment from $A$ to $C$, are glued crosswise; the dashed
arc is one of the two lifts of this segment.  The boundary circles have
widths $4$ and $16$.  The cover is branched over $A$ and $C$, whose
preimages are the ramification points $\widetilde A,\widetilde C$, and $B$ has two
unramified lifts $B_0,B_1$.}
\label{fig:annular-cover}
\end{figure}

The annulus $\overline\Acal$ with its $\widetilde M$-affine structure is thus
determined by the lattice $\widetilde M$ of the first descendent exponents.
Section~\ref{sec:compactification} compares it with the affine bases of
distinguished toric degenerations in the sense of \cite{carletal2024}.

\subsection{Local Kummer charts}
\label{subsec:kummer-charts}

Each matrix in \eqref{eq:lifted-monodromies} exchanges $\alpha$ and $\beta$
in $\widetilde M/M$, while the core matrix $K$ acts trivially.  The
monodromy of the local system with fiber $\widetilde M/M$ therefore acts
through the exchange of $\alpha$ and $\beta$, and $\deg_4$ of
\eqref{eq:diagonal-charge} is the projection onto its largest quotient on
which the monodromy acts trivially.  The local system becomes constant on
the connected double cover
\[
 p_{\mathrm{fr}}:\Acal^\circ_{\mathrm{fr}}\longrightarrow\Acal^\circ
\]
defined by the kernel of the monodromy action on $\widetilde M/M$.  We call
$p_{\mathrm{fr}}$ the \emph{fractional-exponent cover}: it is the smallest
connected cover on which the fractional-exponent group forms a constant
local system.  Its deck involution $\delta_{\mathrm{deck}}$ exchanges the
classes $\alpha$ and $\beta$.  The lifted slab maps require, in addition,
roots of the slab functions; these are the Kummer roots constructed below.

\emph{Placement of the transitions.}  For $V\in\{A,B,C\}$ let $c_V$ be the
outer half of the star edge through $V$, the segment from $V$ to the vertex
$p_V$ of $\Delta$.  The complement of $c_A\cup c_B\cup c_C$ in the interior
of $\Delta$ is simply connected and contains $O$, the inner half-edges and
the branch segment from $A$ to $C$.  Its preimage in $\overline\Acal$
therefore consists of two \emph{sheets} $S_0$ and $S_1$, each mapped
homeomorphically onto it; we write $O_0\in S_0$ and $O_1\in S_1$ for the two
lifts of the central vertex $O$.  The preimage of $c_B$ consists of one arc
from each $B_i$ to the boundary, and the preimages of $c_A$ and $c_C$
consist of two arcs from $\widetilde A$ and two from $\widetilde C$, one to
each boundary circle.  We call these six arcs the \emph{cuts}.  All
transitions at the singular points, namely the affine transitions, the
lifted slab maps below and the line transitions of
Section~\ref{sec:jacobi-theta}, are carried by the cuts, which run from the
singular points outward to the boundary.  The sheets carry no transition;
in particular none is carried by the inner half-edges, by the preimage of
the branch segment, or near $O_0$ and $O_1$.  By
Proposition~\ref{prop:proper-supports}, no wall support of
Section~\ref{sec:jacobi-theta} meets an arc carrying a transition.

We give $S_0$ a fixed frame and $S_1$ the continuation of that frame across the arc
over $c_A$ that $\gamma_A$ crosses from $S_0$, so that this arc carries no
transition.  We place the lifted slab map at each singular
point on a single arc from that point to the boundary: at $B_0$ and $B_1$ on
the arcs over $c_B$, at $\widetilde A$ on the arc over $c_A$ that $\gamma_A$
crosses from $S_1$, and at $\widetilde C$ on the arc over $c_C$ that
$\gamma_C$ crosses from $S_1$.  The arc over $c_C$ that $\gamma_C$ crosses
from $S_0$ carries only the change of affine chart along
$\gamma_{\mathrm{core}}$.  As in Appendix~\ref{app:calculations}, a product
$gh$ of loops traverses $g$ first.  Thus $\gamma_{\mathrm{core}}=\gamma_C\gamma_A^{-1}$
crosses only the arc carrying the change of chart and the arc carrying no
transition, and the transport along it has linear part $K$ and no slab
factor.  The dual graph of this decomposition is the Schreier graph of the
parity subgroup for the transversal $\{1,\gamma_A\}$, and its five edges
other than the transversal edge correspond to the generators
\eqref{eq:parity-matrices}.

\emph{Chart algebras.}  Every vertex of the path tree $\Tcal$ of
Section~\ref{subsec:root-orbit}, a reduced edge path from $S_0$ to a sheet,
carries the frame obtained from that of $S_0$ by transport along the path;
the frames of $S_0$ and $S_1$ fixed above are those at the empty path and at
the path that crosses the arc over $c_A$ carrying no transition.  The
\emph{chart algebra} of a sheet, or of a point or a region of $\Acal^\circ$
read in the frame at a vertex, is $\Bbbk[\widetilde M]$ in that frame, over
the Kummer ring of the relevant finite stage: for a finite rooted subtree
$T\subset\Tcal$ containing the vertex, it is the Kummer ring
$R_T^{\mathrm{root}}$ of \eqref{eq:finite-kummer-ring}, into which
$\Bbbk[\widetilde M]$ in that frame maps by pull-back along the path of the
vertex (Section~\ref{subsec:root-stage-regularity}), tensored with the
coefficient ring in use, such as $R/\mathfrak m^N$ in
\eqref{eq:annular-root-object} or a truncated coefficient ring of
Section~\ref{sec:compactification}.

For $V\in\{\widetilde A,B_0,B_1,\widetilde C\}$ let $u_V\in\widetilde M$ be
the primitive vector spanning the invariant line of $H_V$, and let
$x_V=z^{2\epsilon_Vu_V}$ with the signs $\epsilon_V$ below.  The exponent
$2\epsilon_Vu_V$ is the invariant character $\bar u_V\in M$ of the disk
monodromy $T_V$ about the image of $V$, that is, one of the characters
$\bar u_A,\bar u_B,\bar u_C$ of $T_A,T_B,T_C$ listed in
Appendix~\ref{app:calculations}, transported linearly along $\gamma_A$ for
$B_1$.  The pulled-back slab functions are
\[
 f_{\widetilde A}=1+x_{\widetilde A},\qquad
 f_{B_i}=1+a_\rho x_{B_i}+x_{B_i}^2,\qquad
 f_{\widetilde C}=1+x_{\widetilde C},
\]
of degree $e_V=1,2,2,1$ in $x_V$.  Let $N_V\in\widetilde M^\vee$ be the
primitive conormal of the invariant line, with the sign fixed by
\begin{equation}\label{eq:monodromy-sign}
 H_V-\mathrm{Id}=-\frac{e_V}{d_V}\,(2\epsilon_Vu_V)\,N_V^{T}=-\epsilon_V\,u_VN_V^T,
\end{equation}
where $d_V$ is the root degree below.  The data, with $u_V$ and $N_V$ given
by their $F$-coordinates, are
\begin{equation}\label{eq:root-data}
\begin{array}{c|cccc}
 V&\widetilde A&B_0&B_1&\widetilde C\\ \hline
 F^{-1}u_V&(3,1)&(1,-1)&(5,3)&(1,3)\\
 \epsilon_V&1&1&-1&-1\\
 d_V&2&4&4&2\\
 F^{T}N_V&(1,-3)&(-1,-1)&(-3,5)&(-3,1).
\end{array}
\end{equation}
On the complement of the slab divisor put
\begin{equation}\label{eq:local-root-ring}
 R_V^{\mathrm{root}}
 =\Bbbk[\widetilde M][f_V^{-1},g_V]/(g_V^{d_V}-f_V),
 \qquad
 \widetilde\theta_V(z^m)=g_V^{\langle N_V,m\rangle}z^m.
\end{equation}
The ring $R_V^{\mathrm{root}}$ is the Kummer ring of $f_V$.  Its spectrum is
the restriction, over the complement of the slab divisor, of the Kummer
cover $\Spec\Bbbk[\widetilde M][g_V]/(g_V^{d_V}-f_V)$, whose quotient by
$\mu_{d_V}$, acting by $g_V\mapsto\zeta g_V$ and trivially on characters, is
the $d_V$th root stack of the slab divisor (Appendix~\ref{app:calculations}).
Over the complement of the divisor the Kummer cover is a
$\mu_{d_V}$-torsor, the root stack is the complement itself, and $g_V$ is not
a function on it.  Since $\langle N_V,u_V\rangle=0$, the map
$\widetilde\theta_V$ fixes $f_V$ and $g_V$ and is an automorphism of
$R_V^{\mathrm{root}}$.  It commutes with $\mu_{d_V}$ on $z^m$ exactly when
$d_V$ divides $\langle N_V,m\rangle$.  This holds for $m\in M$, where
$\langle N_V,m\rangle\in4\ZZ$ and $\widetilde\theta_V$ multiplies $z^m$ by an
integral power of $f_V$, and it fails for some $m\in\widetilde M$ because
$N_V$ is primitive.  The lifted slab maps, the line transitions of
Section~\ref{sec:jacobi-theta} and every transport built from them therefore
act on Kummer rings and Kummer covers, and not on root stacks.

The lifted slab map at $V$ is placed on its arc so that the transport along
the loop $\gamma_V$ of Section~\ref{subsec:parity-cover}, based in $S_0$, is
$z^m\mapsto g_V^{\langle N_V,m\rangle}z^{H_Vm}$.  By
\eqref{eq:monodromy-sign}, $z^{H_Vm}=z^{m-\epsilon_V\langle N_V,m\rangle u_V}$,
and $g_V^{d_V}x_V^{-e_V}=f_Vx_V^{-e_V}$ tends to one where $x_V$ is large;
in this sense the transport is trivial to leading order there, which fixes
the sign of $N_V$.

\begin{proposition}\label{prop:minimal-roots}
The formula \eqref{eq:local-root-ring} defines an automorphism of the Kummer
ring $R_V^{\mathrm{root}}$ extending the transition of the disk: on
characters $m\in M$, the transports about $\widetilde A$, $B_0$ and $\widetilde C$
based in $S_0$ are the disk transports about $A$ twice, $B$ once and $C$
twice, and the transport about $B_1$ is the disk transport about $B$ written
in the frame of $S_1$.  For $a_\rho\neq\pm2$ the root degrees $2,4,4,2$ are
minimal.
\end{proposition}

\begin{proof}
The root $g_V$ is a unit on the localized chart and
$\langle N_V,m\rangle$ is integral.  Moreover $N_V=4n_V$, where
$n_V\in M^\vee$ is the primitive conormal of the disk at the image of $V$,
transported linearly along $\gamma_A$ for $B_1$; the arc that $\gamma_A$
crosses from $S_0$ carries no transition.  The disk transport about that
singular point is $z^m\mapsto f_V^{\langle n_V,m\rangle}z^{T_Vm}$, where
$T_V-\mathrm{Id}=-e_V\bar u_Vn_V^T$ by \eqref{eq:disk-rank-one}; at $B_1$, the slab function, the invariant
character, the conormal and the monodromy are those of $B$ transported
linearly along $\gamma_A$, that is, written in the frame of $S_1$.  The disk
transport fixes $f_V$ and $n_V$.  At $\widetilde A$ and $\widetilde C$ the
loop traverses the disk loop twice, so on $m\in M$ the slab factor is
$f_V^{2\langle n_V,m\rangle}=g_V^{\langle N_V,m\rangle}$; at $B_0$ and $B_1$
it traverses it once, giving $f_V^{\langle n_V,m\rangle}=g_V^{\langle
N_V,m\rangle}$.  Since $N_V$ is primitive, $\langle N_V,m\rangle=1$ for some
$m\in\widetilde M$, so the slab factor of that character is a $d_V$th root
of $f_V$.  For $a_\rho\neq\pm2$ the Laurent polynomial $f_V$ vanishes to order
one along each component of its zero divisor, so its class in
$\Bbbk(\widetilde M)^\times/\Bbbk(\widetilde M)^{\times d_V}$ has order
$d_V$, and no proper divisor of $d_V$ suffices.
\end{proof}

The placement above thus defines the lifted coefficient system.  On
characters in $M$ it agrees with the pullback of the disk system germwise at
each singular point (Proposition~\ref{prop:minimal-roots}), but it is not
that pullback.  In the pulled-back system each of the two arcs over $c_A$,
and each of the two over $c_C$, carries the disk transition, whose slab
factor on a single crossing is $f_V^{\langle n_V,m\rangle}$ with
$\langle n_V,m\rangle\in\frac14\ZZ$ for $m\in\widetilde M$.  On characters in
$M$, the pulled-back disk transport along $\gamma_A\gamma_B\gamma_A^{-1}$ is
the transport about $B_1$ conjugated by
$z^m\mapsto f_A^{\langle n_A,m\rangle}z^m$, and the pulled-back transport
along $\gamma_{\mathrm{core}}$ has slab factors, powers of $f_C$ and of the
transport of $f_A$, whereas the transport of the lifted system along
$\gamma_{\mathrm{core}}$ is the monomial map with linear part $K$.

At the two cusps $a_\rho=\pm2$ the degree at $B_0$ and $B_1$ drops to two.  If
$f_B(a_\rho,z)=1+a_\rho z+z^2$, then
\begin{equation}\label{eq:cusp-roots}
\begin{aligned}
 a_\rho=2:&\quad g^4-f_B=(g^2-(1+z))(g^2+(1+z)),\\
 a_\rho=-2:&\quad g^4-f_B=(g^2-(1-z))(g^2+(1-z)).
\end{aligned}
\end{equation}
Away from the reduced zero divisor these are two connected quadratic
torsors; their closures meet over the divisor.  We retain this fourth-root
stack rather than replacing it by one connected square-root cover.

\subsection{The infinite root orbit}
\label{subsec:root-orbit}

The failure of finite Kummer closure already occurs on the overlap of the
$\widetilde A$ and $B_0$ charts.  Write
\[
 x_A=x_{\widetilde A}=z^{2u_{\widetilde A}},\qquad
 x_B=x_{B_0}=z^{2u_{B_0}},\qquad f_A=1+x_A,\qquad
 f_B(x)=1+a_\rho x+x^2,
\]
and let $\mathcal K=\Bbbk(x_A,x_B)$.  Since
$\langle N_{\widetilde A},2u_{B_0}\rangle=8$ and $d_{\widetilde A}=2$, the
$\widetilde A$-slab shear fixes $f_A$ and sends
\begin{equation}\label{eq:shear}
 x_B\longmapsto f_A^4x_B.
\end{equation}
The transport about $\widetilde A$ also applies the linear monodromy, and
$H_A(2u_{B_0})=2u_{B_0}-4\cdot2u_{\widetilde A}$; it therefore fixes $x_A$
and sends $x_B$ to $x_A^{-4}f_A^4x_B=(1+x_A^{-1})^4x_B$.  For $a_\rho\neq\pm2$
write $f_B(x)=(1+\lambda x)(1+\lambda^{-1}x)$ with $\lambda+\lambda^{-1}=a_\rho$
and $\lambda^2\neq1$.

\begin{theorem}[Infinite Kummer orbit]\label{thm:infinite-root-orbit}
Let $u\in\Bbbk(x_A)$ be nonconstant, and let $\omega_u$ be the automorphism
of $\mathcal K$ fixing $x_A$ with $\omega_u(x_B)=ux_B$.  The slab shear
\eqref{eq:shear}, with $u=f_A^4$, and the transport about $\widetilde A$,
with $u=(1+x_A^{-1})^4$, are of this form.
\begin{enumerate}
\item If $a_\rho\neq\pm2$, closure of a fourth root of $f_B(x_B)$ under all
powers of $\omega_u$ requires the fourth roots $g_n$ of
\begin{equation}\label{eq:root-tower}
 g_n^4=f_B(u^nx_B)
 =\bigl(1+\lambda u^nx_B\bigr)\bigl(1+\lambda^{-1}u^nx_B\bigr),
 \qquad n\in\ZZ,
\end{equation}
and the classes of these radicands are linearly independent over $\ZZ/4$ in
$\mathcal K^\times/\mathcal K^{\times4}$.
\item If $a_\rho=\pm2$, then $f_B(u^nx_B)=(1\pm u^nx_B)^2$, and the classes of
$1\pm u^nx_B$, $n\in\ZZ$, are linearly independent over $\ZZ/2$ in
$\mathcal K^\times/\mathcal K^{\times2}$.
\item In either case no finite extension of $\mathcal K$ containing a fourth root
of $f_B(x_B)$ admits an automorphism extending $\omega_u$.  In particular no
finite algebraic Kummer extension containing the initial $\widetilde A$- and
$B_0$-roots is stable under the slab shear or under transport about
$\widetilde A$.
\end{enumerate}
\end{theorem}

\begin{proof}
Suppose $a_\rho\neq\pm2$.  In the principal ideal domain $\Bbbk(x_A)[x_B]$,
\[
 f_B(u^nx_B)
 =u^{2n}\bigl(x_B+\lambda^{-1}u^{-n}\bigr)\bigl(x_B+\lambda u^{-n}\bigr).
\]
The elements $\lambda^{\pm1}u^{-n}$, $n\in\ZZ$, of $\Bbbk(x_A)$ are pairwise
distinct: an equality $\lambda^{\epsilon}u^{-n}=\lambda^{\epsilon'}u^{-m}$
makes $u^{m-n}$ constant, hence $m=n$ because $u$ is nonconstant, and then
$\lambda^{\epsilon}=\lambda^{\epsilon'}$ forces $\epsilon=\epsilon'$ because
$\lambda^2\neq1$.  Thus $\pi_n=x_B+\lambda^{-1}u^{-n}$ is a prime dividing the
$n$th radicand exactly once and no other radicand.  If a finite product
$\prod_nf_B(u^nx_B)^{k_n}$ is a fourth power in $\mathcal K$, its $\pi_n$-adic
valuation $k_n$ is divisible by four.  This proves (1).  For (2), the same
argument applies to the primes $x_B\pm u^{-n}$ of $1\pm u^nx_B$.

For (3), let $\mathcal K'\supset\mathcal K$ be finite, let $\widetilde\omega$ be an
automorphism of $\mathcal K'$ extending $\omega_u$, and let $h\in\mathcal K'$ satisfy
$h^4=f_B(x_B)$.  Then $\widetilde\omega^n(h)^4=f_B(u^nx_B)$ for every
$n\in\ZZ$.  If $a_\rho\neq\pm2$, the field $\Bbbk$ contains the fourth roots of
unity, so by Kummer theory and (1) the subfield generated by the
$\widetilde\omega^n(h)$ with $|n|\leq r$ has degree $4^{2r+1}$ over $\mathcal K$
for every $r$, which is impossible.  If $a_\rho=\pm2$, then
$\widetilde\omega^n(h)^2=\pm(1\pm u^nx_B)$, and since $-1$ is a square in
$\Bbbk$, $\mathcal K'$ contains square roots of all $1\pm u^nx_B$; by (2) the same
degree argument applies with $2^{2r+1}$.
\end{proof}

\emph{The path tree.}  Let $\mathfrak C$ be the graph whose vertices are the
two sheets $S_0,S_1$ and whose edges are the six cuts, each joining the
sheets on its two sides.  The sheets are simply connected and the cuts are
contractible, so the fundamental group of $\mathfrak C$ based at $S_0$ is
$\Pi=\pi_1(\Acal^\circ)$, a free group of rank five.  The \emph{path tree}
$\Tcal$ is the universal covering tree of the graph $\mathfrak C$.  Its
vertices are the reduced edge paths in $\mathfrak C$ beginning at $S_0$,
equivalently the homotopy classes of paths in $\Acal^\circ$ from a base
point $b_0\in S_0$ to $b_0$ or to the other lift $b_1\in S_1$ of its image
in $\Delta$; its root is the empty path at $S_0$.

Let $\mathbb T=\Spec\Bbbk[\widetilde M]$ be the torus of the sheet $S_0$, and let
$T\subset\Tcal$ be a finite rooted subtree.  We call an edge of $T$ a
\emph{Kummer edge} if it crosses, in either direction, an arc carrying a
lifted slab map; edges crossing the other two arcs adjoin no root.
Number the Kummer edges $e=1,\dots,\ell$, the edge $e$ crossing the arc of the
singular point $V_e$.  It has the root degree $d_e=d_{V_e}$ and the
\emph{transported radicand} $\varrho_e$, the transport of $f_{V_e}$ to $S_0$
along the path of $T$ preceding $e$.  The vectors $u_V$ lie in the
sublattice of $\widetilde M$ of elements whose $F$-coordinates $(m_1,m_2)$
satisfy $m_1\equiv m_2\pmod2$, which the matrices
\eqref{eq:parity-matrices} preserve and on which every $N_V$ is even.  A
lifted slab map therefore multiplies a transported character
$z^{2\epsilon_{V'}u_{V'}}$ by a power of the root it uses with exponent
divisible by four, that is, by an integral power of the corresponding
transported radicand.  By induction along the path, every $\varrho_e$ is a
rational function on $\mathbb T$.  Let $U_T\subset\mathbb T$ be the complement of the zeros
and poles of $\varrho_1,\dots,\varrho_\ell$, and put
$G_T=\prod_e\mu_{d_e}$.  The \emph{Kummer ring} of $T$ is
\begin{equation}\label{eq:finite-kummer-ring}
 R_T^{\mathrm{root}}
 =\mathcal O(U_T)[r_1,\dots,r_\ell]\big/\bigl(r_e^{d_e}-\varrho_e\bigr)_{e},
\end{equation}
on which $\zeta\in G_T$ acts by $r_e\mapsto\zeta_er_e$.  It is finite étale
over $\mathcal O(U_T)$, and $\Spec R_T^{\mathrm{root}}$ is a $G_T$-torsor
over $U_T$.  The lifted slab map at the Kummer edge $e$ uses the root $r_e$,
the transport of $g_{V_e}$, so transport along the paths of $T$ is defined
on $R_T^{\mathrm{root}}$.  Let $D_e$ be the divisor of zeros of
$\varrho_e$ on $\mathbb T$.  The \emph{finite root stack} of $T$ is
\[
 \Xfrak_T=\sqrt[d_1]{(\mathbb T,D_1)}\times_{\mathbb T}\cdots\times_{\mathbb T}\sqrt[d_\ell]{(\mathbb T,D_\ell)},
\]
in the convention \eqref{eq:root-stack-chart-app}.  Over $U_T$ it is $U_T$
itself, the quotient of $\Spec R_T^{\mathrm{root}}$ by $G_T$.
Section~\ref{sec:compactification} extends $\Spec R_T^{\mathrm{root}}$ to a
Kummer cover $P_T$ of $\mathbb T$ with $\Xfrak_T=[P_T/G_T]$; we call $\Xfrak_T$ with
its atlas $P_T$ the \emph{finite root-stack presentation} of $T$.
Transport acts on $R_T^{\mathrm{root}}$, that is, on $P_T$ over $U_T$, and
not on $\Xfrak_T$, because the lifted slab maps do not commute with $G_T$.
By Appendix~\ref{app:calculations}, the stabilizer of $\Xfrak_T$ at a
geometric point over $z\in\mathbb T$ is the product of the factors
$\mu_{d_e}$ with $z\in D_e$.  If $T\subset T'$, localization
followed by adjunction of the new roots maps $R_T^{\mathrm{root}}$
injectively to $R_{T'}^{\mathrm{root}}$, and forgetting the new roots gives
$\Xfrak_{T'}\to\Xfrak_T$.  The paths $\gamma_{\widetilde A}^n$ followed by a
crossing of the $B_0$-arc give the radicands $f_B(u^nx_B)$,
$u=(1+x_A^{-1})^4$, of Theorem~\ref{thm:infinite-root-orbit}.

\begin{definition}\label{def:pathwise-root-object}
Let $R=\Bbbk[t_1,t_2]$ be the deformation coefficient algebra, where
$t_1,t_2$ are the singleton parameters, and let $\mathfrak m=(t_1,t_2)$ be the
singleton ideal.  The \emph{pathwise Kummer ring} is
$R_{\mathrm{path}}^{\mathrm{root}}=\varinjlim_{T\subset\Tcal}R_T^{\mathrm{root}}$,
the colimit over the finite rooted subtrees $T\subset\Tcal$.
The pathwise root object and its adic completion are the formal systems
\begin{equation}\label{eq:pathwise-root-object}
 \Xfrak_{\mathrm{path}}=\varprojlim_{T\subset\Tcal}\Xfrak_T,
 \qquad
 \widehat\Xfrak_{\mathrm{path}}
 =\varinjlim_N\varprojlim_{T\subset\Tcal}
  \bigl(\Xfrak_T\times\Spec R/\mathfrak m^N\bigr).
\end{equation}
The group $\Pi$ acts on $\Tcal$ by deck transformations, a loop acting by
concatenation, and through transport on $R_{\mathrm{path}}^{\mathrm{root}}$,
the Kummer root of an edge going to the Kummer root of its translate.  The
\emph{annular coefficient system} is $\widehat\Xfrak_{\mathrm{path}}$
together with this action of $\Pi$.  A structure over it, such as a module,
an algebra, a bracket or a wall structure, is a $\Pi$-equivariant compatible
system of such structures over the rings
\begin{equation}\label{eq:annular-root-object}
 R_T^{\mathrm{root}}\otimes_\Bbbk R/\mathfrak m^N,
 \qquad T\subset\Tcal\ \text{a finite rooted subtree},\quad N\geq1.
\end{equation}
The coefficient ring at every vertex of $\Tcal$ is
$R_{\mathrm{path}}^{\mathrm{root}}\otimes_\Bbbk R$, transport along the path
of the vertex identifying it with the ring at the empty path.  A
\emph{finite transport diagram modulo $\mathfrak m^N$} is a finite diagram of
finitely presented modules or algebras over the coefficient rings of
finitely many vertices of $\Tcal$, reduced modulo $\mathfrak m^N$, together with
finitely many homomorphisms, transports along edges of $\Tcal$, and
identities among their composites.
\end{definition}

The limit over $T$ is formed along the root-forgetting maps, the colimit of
Kummer rings along the injections $R_T^{\mathrm{root}}\to R_{T'}^{\mathrm{root}}$,
and the colimit over $N$ along the closed immersions
$\Spec R/\mathfrak m^N\to\Spec R/\mathfrak m^{N+1}$.  The limits are formal systems: by
Theorem~\ref{thm:forced-category}, $\Xfrak_{\mathrm{path}}$ is not
represented by a formal scheme or by an algebraic stack whose diagonal is
locally of finite type.  The group $\Pi$ acts on the Kummer rings,
preserving the root degrees, but not on the root stacks $\Xfrak_T$, whose
root groups $G_T$ do not commute with the lifted slab maps.  For this reason
no quotient by $\Pi$ is formed: structures over the annular coefficient
system are defined through the Kummer rings \eqref{eq:annular-root-object}
and the action of $\Pi$ on them.  This inverse system is of the same
formal kind as the infinite root stacks of \cite{talpovistoli2018}; the
additional feature here is that the roots are indexed by reduced transport
paths and permuted by the deck group $\Pi$.  Wall structures over this
object and their consistency are the subject of
Section~\ref{sec:jacobi-theta}.

\begin{theorem}[Finite-stage definability and nonrepresentability]
\label{thm:forced-category}
Every finite transport diagram modulo $\mathfrak m^N$ is defined over
$R_T^{\mathrm{root}}\otimes_\Bbbk R/\mathfrak m^N$ for some finite rooted subtree $T$,
that is, over $U_T$ on the Kummer cover $P_T$ of the finite root-stack
presentation of $T$.  Any two such subtrees have a common finite refinement.
On the other hand, if all transported roots and their automorphisms are
retained, the pathwise root object $\Xfrak_{\mathrm{path}}$ is not
represented by a formal scheme, a finite root stack, or an algebraic stack
whose diagonal is locally of finite type.
\end{theorem}

\begin{proof}
A finite transport diagram involves finitely many vertices of $\Tcal$ and,
modulo $\mathfrak m^N$, finitely many deformation bidegrees.  The rooted union of the
corresponding paths is a finite subtree.  The ring
$R_{\mathrm{path}}^{\mathrm{root}}$ is the colimit of the Kummer rings of the
finite subtrees (Appendix~\ref{app:calculations}), so by finite
presentation the finitely many generators, relations, arrows and identities
are defined over one of them after a further finite enlargement.  Rooted union gives a common
refinement.

For the negative assertion, let $u=(1+x_A^{-1})^4$, let $\lambda$ be a root
of $\lambda^2-a_\rho\lambda+1$ (so $\lambda=\pm1$ when $a_\rho=\pm2$), and let $w$ be a
geometric point of $\mathbb T$ with $x_A=-\frac12$ and $x_B=-\lambda^{-1}$.  There
$f_A=\frac12$ and $u=1$, so every radicand $f_B(u^nx_B)$ in
\eqref{eq:root-tower} vanishes.  For a finite set $S\subset\ZZ$, let $T$ be
the rooted union of the paths $\gamma_{\widetilde A}^n$ followed by a
crossing of the $B_0$-arc, $n\in S$.  Since the transport about
$\widetilde A$ fixes $f_A$, every root adjoined along
$\gamma_{\widetilde A}^n$ has the radicand $f_A$, which does not vanish at
$w$.  The root adjoined at the final crossing has the radicand
$\varrho_n=f_B(u^nx_B)$, which is regular at $w$, because $x_A$ and $f_A$ are units
there, and vanishes at $w$; hence $w\in D_n$.  The stabilizer of
$\Xfrak_T$ at the point over $w$ is therefore $\prod_{n\in S}\mu_4$.
The automorphism group of a compatible point of $\Xfrak_{\mathrm{path}}$ over
$w$ is the inverse limit of the stabilizers over all finite subtrees, and it
contains
\[
 \prod_{n\in\ZZ}\mu_4
\]
as a closed subgroup.  This group scheme is not of finite type.  A formal
scheme has trivial stabilizers, a finite root stack has finite-type
diagonal, and an algebraic stack with diagonal locally of finite type cannot
have such a geometric stabilizer.
\end{proof}

Theorem~\ref{thm:forced-category} is the only property of the inverse
system used later: every product, period and logarithmic chart below is
constructed over the Kummer ring of a finite stage, and the pro-object
records the compatibility of these constructions.
